\documentclass[a4paper, 12p]{article}
\usepackage[spanish]{babel} 
\usepackage{amsmath} 
\usepackage[colorlinks=true]{hyperref}
\usepackage{enumitem} 
\usepackage{graphicx}   
\usepackage[a4paper,top=3cm,bottom=3cm,left=3cm,right=3cm,marginparwidth=1.75cm]{geometry} 
\usepackage[]{subfigure}
\graphicspath{{./graficos/Lab-Termo/Git/Documentos}} 
\usepackage{float}
%\usepackage[]{txfonts}
\usepackage{multicol}



\newenvironment{Figura}
{\par\medskip\noindent\minipage{\linewidth}}
{\endminipage\par\medskip}

%==================================================================================
\begin{document}
\begin{titlepage}
      \begin{center}     
              
            \includegraphics[width=0.2\textwidth]{img/escudo_udec.png}                       %Para poner logo udec   %{nombre carpeta\nombreimagen}
            
            
            
            \vspace{1cm}
            \textsc{{\LARGE Universidad de Concepción}}
            
            \vspace{1cm}
            {\scshape\Large Facultad de ciencias fisícas y matemáticas \par}
            \vspace{2cm}
            {\scshape\Huge Laboratorio 4 \par}
            \vspace{2cm}
            {\itshape\Large Proyecto laboratorio termodinámica \par}
            \vfill
            {\Large Autores: \par}
            {\Large Martina Contreras, Noemí De La Peña, Benjamín Opazo. \par}
            \vfill
            \vfill
            {\Large Profesor: \par}
            {\Large Claudio Alonso Faúndez Araya \par}
            \vfill
            \vfill
            {\Large Carrera: \par}
            {\Large Ciencias fisícas \par}
            \vfill
            \vfill
            {\Large Ayudante: \par}
            {\Large Renata Hernández Lopez \par}
            \vfill
            {\Large Noviembre 2022 \par}
      \end{center}
\end{titlepage}            
 

\tableofcontents
\newpage

%========Introducción
\section{Introducción}
En este informe presentaremos la experiencia del laboratorio, cuya finalidad fue la calcular la densidad de aceite de cocina por medio del ley de la hidrostática. primero presentaremos los objetivos del laboratorio, para luego definir conceptos importantes, tales como, que es la ley de la hidrostática y la densidad. Luego, responderemos las preguntas propuestas en el análisis. Para finalizar, anunciaremos si los objetivos se lograron cumplir junto con la veracidad de nuestra hipótesis.

\section{Objetivos}
\begin{itemize}
      \item Determinar la densidad de aceite de cocina por medio de la ley de la Hidrostática.
      
      \item Comparar el valor determinado con el de literatura.
\end{itemize}



%=========Marco Teórico
\section{Marco Teórico} 



%=========Materiales
\section{Materiales}


\begin{itemize}
      \item 1 tubo en U.
      \item 1 soporte universal.
      \item 2 pinzas de sujeción.
      \item Regla escuadra.
      \item Aceite de cocina.
      \item agua,
\end{itemize}

\section{Procedimiento}

\begin{itemize}
      \item Primero debemos llenar el tubo en u con $8[cc]$ de agua.
      
      \item Luego agregamos $4[cc]$ de aceite de cocina por algunas de las ramas del tubo en u.
      
      \item Medir las alturas mostradas en la Fig.1 para calcular la altura de cada una de las columnas.
      
\end{itemize}

\section{Análisis}




% \begin{thebibliography}{6}
%       \bibitem{libro} Sears, F. W., Salinger, G. L. \& Peris, A. J. (2021, 10 enero). Termodinámica, teoría cinética y termodinámica estadística (Spanish Edition) (1.a ed.). Reverte.
%       \bibitem{profe} Faundez Araya, C. A. (s. f.). Capítulo 3. Primer Principio de la Termodinámica [Diapositivas]. Universidad de Concepción, 2022.
%       \bibitem{agua}Fundación Aquae. (2021, 10 agosto). CP del agua: significado y valores - Fundación Aquae. FundaciÃ3n Aquae. 
%       \url{https://www.fundacionaquae.org/wiki/cp-del-agua/}
%       \bibitem{acero}DeterminaciÃ3n del calor específico de un sÃ3lido. (s. f.). Recuperado 25 de octubre de 2022,
%       de \url{http://www.sc.ehu.es/sbweb/fisica/estadistica/otros/calorimetro/calorimetro.htm}
% \end{thebibliography}













\end{document}
